\documentclass[12pt, a4paper]{article}
\usepackage[utf8]{inputenc}
\usepackage{amsmath, amssymb, amsthm, amsfonts}
\usepackage{CJKutf8}
\usepackage[margin=1in]{geometry}
\usepackage{tikz}
\newtheorem{lemma}{Lemma}
\newcommand{\g}{\mathfrak{g}}
\newcommand{\h}{\mathfrak{h}}
\newcommand{\R}{\mathbb{R}}
% 重新定义 solution 环境，保留标识并在末尾留出 8cm 空白
\newenvironment{solution}
  {\paragraph{\textit{Solution:}}\par\medskip}
  {\hfill$\blacksquare$\vspace{0.1cm}}

\title{Lie Theory}
\author{
    \begin{CJK*}{UTF8}{gbsn}
    钟皓宇 (12533556)
    \end{CJK*}
}
\date{\today}

\begin{document}

\maketitle
\section{Basic Concepts}

\subsection{ideal and homomorphism}

\subsubsection{Problem 1}

Prove that the set of all inner derivations \(\operatorname{ad} x\), \(x \in L\),
is an ideal of \(\operatorname{Der} L\).

\paragraph{Hint.}
Use the derivation identity. For \(D \in \operatorname{Der}L\), compute
\([D,\operatorname{ad}x](y)\) and show that it equals \(\operatorname{ad}(Dx)(y)\).


Recall that a \emph{derivation} of a Lie algebra \(L\) is a linear map
\(D:L\to L\) satisfying
\[
D([x,y])=[D(x),y]+[x,D(y)]
\]
for all \(x,y\in L\). The set of all derivations is denoted by
\(\operatorname{Der}L\), and it is a Lie algebra under the commutator bracket
\([D_1,D_2]=D_1D_2-D_2D_1\). For each \(x\in L\), the map
\[
\operatorname{ad}x:L\to L,\qquad y\mapsto [x,y],
\]
is called an \emph{inner derivation}.
\subsubsection{Problem 2}

Show that \(\mathfrak{sl}(n,F)\) is precisely the derived algebra of
\(\mathfrak{gl}(n,F)\). Cf. Exercise 1.9.

\paragraph{Hint.}
First observe that every commutator has trace zero. Conversely, use the matrix
units \(E_{ij}\). For \(i \neq j\), get \(E_{ij}\) as a commutator, and get
\(E_{ii}-E_{jj}\) from \([E_{ij},E_{ji}]\).

\subsubsection{Problem 3}

Prove that the center of \(\mathfrak{gl}(n,F)\) equals \(\mathfrak{s}(n,F)\),
the scalar matrices. Prove that \(\mathfrak{sl}(n,F)\) has center \(0\), unless
\(\operatorname{char} F\) divides \(n\), in which case the center is
\(\mathfrak{s}(n,F)\).

\paragraph{Hint.}
Let \(A=(a_{ij})\) commute with every matrix unit \(E_{ij}\). This forces
\(A\) to be scalar. Then intersect the scalar matrices with
\(\mathfrak{sl}(n,F)\): a scalar matrix \(\lambda I\) lies in
\(\mathfrak{sl}(n,F)\) exactly when \(n\lambda=0\).

\subsubsection{Problem 4}

Show that, up to isomorphism, there is a unique Lie algebra over \(F\) of
dimension \(3\) whose derived algebra has dimension \(1\) and lies in \(Z(L)\).

\paragraph{Hint.}
Choose \(0 \neq z \in [L,L]\). Since \([L,L]\subset Z(L)\), choose
\(x,y\in L\) whose images form a basis of \(L/Fz\). After rescaling, arrange
\([x,y]=z\). This gives the three-dimensional Heisenberg Lie algebra.

\subsubsection{Problem 5}

Suppose \(\dim L=3\) and \(L=[L,L]\). Prove that \(L\) must be simple. Observe
first that any homomorphic image of \(L\) also equals its derived algebra.
Recover the simplicity of \(\mathfrak{sl}(2,F)\), where
\(\operatorname{char} F \neq 2\).

\paragraph{Hint.}
If \(I\) is a nonzero proper ideal, then \(L/I\) has dimension \(1\) or \(2\).
But \(L/I\) is also perfect, because \(L=[L,L]\). Show that no nonzero Lie
algebra of dimension \(1\) or \(2\) is perfect. For
\(\mathfrak{sl}(2,F)\), check directly that \(x,y,h\) all lie in the derived
algebra.

\subsubsection{Problem 6}

Prove that \(\mathfrak{sl}(3,F)\) is simple, unless \(\operatorname{char} F=3\).
Cf. Exercise 3. Use the standard basis \(h_1,h_2,e_{ij}\), where \(i \neq j\).
If \(I \neq 0\) is an ideal, then \(I\) is the direct sum of eigenspaces for
\(\operatorname{ad} h_1\) or \(\operatorname{ad} h_2\); compare the eigenvalues
of \(\operatorname{ad} h_1\) and \(\operatorname{ad} h_2\) acting on the
\(e_{ij}\).

\paragraph{Hint.}
Use the root-space decomposition with respect to the diagonal elements
\(h_1,h_2\). A nonzero ideal is stable under \(\operatorname{ad}h_1\) and
\(\operatorname{ad}h_2\), so it contains some root vector \(e_{ij}\) or a
nonzero diagonal element. From one root vector, use brackets to generate all
root vectors and the \(h_i\). If \(\operatorname{char}F=3\), the scalar matrix
\(I_3\) lies in \(\mathfrak{sl}(3,F)\), giving a nonzero center.

\subsubsection{Problem 7}
Let \(L\) be a Lie algebra and let \(A\subset L\) be a subalgebra. The
\emph{normalizer} of \(A\) in \(L\) is
\[
N_L(A):=\{x\in L\mid [x,A]\subset A\}.
\]
Equivalently, \(x\in N_L(A)\) means that taking Lie brackets with \(x\)
preserves the subalgebra \(A\). We say that \(A\) is \emph{self-normalizing} if
\[
N_L(A)=A.
\]


Prove that \(\mathfrak{t}(n,F)\) and \(\mathfrak{b}(n,F)\) are
self-normalizing subalgebras of \(\mathfrak{gl}(n,F)\), whereas
\(\mathfrak{n}(n,F)\) has normalizer \(\mathfrak{t}(n,F)\).

\paragraph{Hint.}
Write \(X=(x_{ij})\in \mathfrak{gl}(n,F)\) and impose the condition
\([X,A]\in \mathfrak{t}(n,F)\), \([X,A]\in \mathfrak{b}(n,F)\), or
\([X,A]\in \mathfrak{n}(n,F)\) for all \(A\) in the relevant subalgebra. Test
this condition on suitable matrix units \(E_{ij}\). The entries of \(X\) below
or above the diagonal are forced to vanish according to the triangular shape.

\subsubsection{Problem 8}

Prove that in each classical linear Lie algebra (1.2), the set of diagonal
matrices is a self-normalizing subalgebra, when \(\operatorname{char} F=0\).

\paragraph{Hint.}
Let \(\mathfrak{h}\) be the diagonal subalgebra. If \(X\) normalizes
\(\mathfrak{h}\), then \([X,H]\in \mathfrak{h}\) for every \(H\in\mathfrak{h}\).
But the off-diagonal entries of \([X,H]\) are multiples of the off-diagonal
entries of \(X\). In characteristic \(0\), choose \(H\) separating the relevant
weights; this forces all off-diagonal entries of \(X\) to vanish.

\subsubsection{Problem 9}

Prove Proposition 2.2.

\paragraph{Hint.}
Check each assertion directly from the definitions. For quotient statements,
first prove that the bracket is well-defined modulo the ideal. For isomorphism
statements, construct the natural homomorphism and identify its kernel and
image.

\subsubsection{Problem 10}

Let \(\sigma\) be the automorphism of \(\mathfrak{sl}(2,F)\) defined in (2.3).
Verify that
\[
\sigma(x)=-y, \qquad \sigma(y)=-x, \qquad \sigma(h)=-h .
\]

\paragraph{Hint.}
Use the explicit matrix formula for \(\sigma\) from (2.3). Substitute the
standard matrices \(x,y,h\) and multiply directly.

\subsubsection{Problem 11}

If \(L=\mathfrak{sl}(n,F)\) and \(g \in \operatorname{GL}(n,F)\), prove that the
map from \(L\) to itself defined by
\[
x \longmapsto -g x^t g^{-1},
\]
where \(x^t\) denotes the transpose of \(x\), belongs to \(\operatorname{Aut} L\).
When \(n=2\) and \(g\) is the identity matrix, prove that this automorphism is
inner.

\paragraph{Hint.}
Use the identity
\[
[x,y]^t=-[x^t,y^t].
\]
This shows that the extra minus sign turns transpose into a Lie algebra
automorphism. For \(n=2\), show that
\[
-x^t=JxJ^{-1},
\qquad
J=
\begin{pmatrix}
0&1\\
-1&0
\end{pmatrix}.
\]

\subsubsection{Problem 12}

Let \(L\) be an orthogonal Lie algebra of type \(B_\ell\) or \(D_\ell\). If
\(g\) is an orthogonal matrix, in the sense that \(g\) is invertible and
\(g^t s g=s\), prove that
\[
x \longmapsto gxg^{-1}
\]
defines an automorphism of \(L\).

\paragraph{Hint.}
Write the orthogonal Lie algebra as
\[
L=\{x\in \mathfrak{gl}(n,F)\mid x^t s+sx=0\}.
\]
For \(x\in L\), compute
\[
(gxg^{-1})^t s+s(gxg^{-1})
\]
using \(g^t s g=s\). This proves that conjugation by \(g\) preserves \(L\).
Conjugation automatically preserves the Lie bracket.
\subsection{Solvable and Nilpotent Lie Algebras}

\subsubsection{Problem 1}

Let \(I\) be an ideal of \(L\). Then each member of the derived series or
descending central series of \(I\) is also an ideal of \(L\).

\paragraph{Hint.}
For the derived series, use induction and the Jacobi identity. If \(I^{(r)}\)
is an ideal of \(L\), then
\[
I^{(r+1)}=[I^{(r)},I^{(r)}].
\]
To prove \(I^{(r+1)}\) is an ideal, compute \([x,[a,b]]\) with
\(x\in L\) and \(a,b\in I^{(r)}\). For the descending central series, use the
same idea with \(\gamma_{r+1}(I)=[I,\gamma_r(I)]\).

\subsubsection{Problem 2}

Prove that \(L\) is solvable if and only if there exists a chain of subalgebras
\[
L=L_0 \supset L_1 \supset L_2 \supset \cdots \supset L_k=0
\]
such that \(L_{i+1}\) is an ideal of \(L_i\) and such that each quotient
\(L_i/L_{i+1}\) is abelian.

\paragraph{Hint.}
If \(L\) is solvable, take the derived series
\[
L=L^{(0)}\supset L^{(1)}\supset L^{(2)}\supset \cdots .
\]
The quotient \(L^{(i)}/L^{(i+1)}\) is abelian by definition. Conversely, use
induction on the length of the chain and the fact that an extension of a
solvable Lie algebra by an abelian Lie algebra is solvable.

\subsubsection{Problem 3}

Let \(\operatorname{char}F=2\). Prove that \(\mathfrak{sl}(2,F)\) is
nilpotent.

\paragraph{Hint.}
Use the standard basis \(x,y,h\) of \(\mathfrak{sl}(2,F)\). In characteristic
\(2\),
\[
[h,x]=2x=0,\qquad [h,y]=-2y=0.
\]
Also \([x,y]=h\). Hence \(h\) is central and
\[
[\mathfrak{sl}(2,F),\mathfrak{sl}(2,F)]=Fh.
\]
Then the next term of the descending central series is zero.

\subsubsection{Problem 4}

Prove that \(L\) is solvable, respectively nilpotent, if and only if
\(\operatorname{ad}L\) is solvable, respectively nilpotent.

\paragraph{Hint.}
Use the natural Lie algebra homomorphism
\[
\operatorname{ad}:L\to \operatorname{Der}L.
\]
Its kernel is \(Z(L)\), so
\[
\operatorname{ad}L\cong L/Z(L).
\]
For solvability, use that central extensions of solvable Lie algebras are
solvable. For nilpotency, use that if \(L/Z(L)\) is nilpotent, then \(L\) is
nilpotent, with nilpotency class at most one larger.

\subsubsection{Problem 5}

Prove that the nonabelian two dimensional algebra constructed in (1.4) is
solvable but not nilpotent. Do the same for the algebra in Exercise 1.2.

\paragraph{Hint.}
For the two-dimensional nonabelian algebra, choose a basis \(x,y\) with
\[
[x,y]=y.
\]
Then
\[
[L,L]=Fy,
\qquad
[[L,L],[L,L]]=0,
\]
so \(L\) is solvable. But the descending central series satisfies
\[
L^2=[L,L]=Fy,\qquad L^3=[L,L^2]=Fy,
\]
and therefore never becomes zero. Hence \(L\) is not nilpotent. For the algebra
in Exercise 1.2, compute its derived series and descending central series in
the same way.

\subsubsection{Problem 6}

Prove that the sum of two nilpotent ideals of a Lie algebra \(L\) is again a
nilpotent ideal. Therefore, \(L\) possesses a unique maximal nilpotent ideal.
Determine this ideal for each algebra in Exercise 5.

\paragraph{Hint.}
Let \(I,J\) be nilpotent ideals. First note that \(I+J\) is again an ideal.
Use the fact that
\[
[I,J]\subset I\cap J.
\]
One way to finish is to use Engel's theorem: show that every element of
\(I+J\) acts nilpotently under the adjoint representation. The sum of all
nilpotent ideals is then itself nilpotent, hence gives the unique maximal
nilpotent ideal. This ideal is usually called the nilradical.

\subsubsection{Problem 7}

Let \(L\) be nilpotent, and let \(K\) be a proper subalgebra of \(L\). Prove
that \(N_L(K)\) includes \(K\) properly.

\paragraph{Hint.}
Consider the action of \(K\) on the quotient vector space \(L/K\) by
\[
k\cdot (x+K)=[k,x]+K.
\]
Since \(L\) is nilpotent, the operators \(\operatorname{ad}k\) are nilpotent.
By Engel's theorem, there is a nonzero vector \(x+K\in L/K\) killed by \(K\).
This means
\[
[K,x]\subset K.
\]
Thus \(x\in N_L(K)\), but \(x\notin K\).

\subsubsection{Problem 8}

Let \(L\) be nilpotent. Prove that \(L\) has an ideal of codimension \(1\).

\paragraph{Hint.}
Use induction on \(\dim L\). Since \(L\) is nilpotent, its center \(Z(L)\) is
nonzero. Choose \(0\neq z\in Z(L)\), and pass to the quotient \(L/Fz\). By
induction, \(L/Fz\) has an ideal of codimension \(1\). Take its inverse image
in \(L\).

\subsubsection{Problem 9}

Prove that every nilpotent Lie algebra \(L\) has an outer derivation; see
(1.3). Proceed as follows. Write
\[
L=K+Fx
\]
for some ideal \(K\) of codimension one, using Exercise 8. Then
\[
C_L(K)\neq 0.
\]
Choose \(n\) so that
\[
C_L(K)\not\subseteq L^n,
\qquad
C_L(K)\subseteq L^{n-1},
\]
and let
\[
z\in C_L(K)-L^n.
\]
Then the linear map \(\delta\) sending \(K\) to \(0\) and \(x\) to \(z\) is an
outer derivation.

\paragraph{Hint.}
First, \(C_L(K)\neq 0\) because \(K\) is nilpotent, so \(Z(K)\neq 0\), and
\(Z(K)\subset C_L(K)\). To prove that \(\delta\) is a derivation, check the
Leibniz rule separately on brackets of elements in \(K\), and on brackets
\([x,k]\) with \(k\in K\). The condition \(z\in C_L(K)\) is exactly what is
needed. To prove that \(\delta\) is outer, suppose \(\delta=\operatorname{ad}y\).
Since \(\delta(K)=0\), one gets \(y\in C_L(K)\). Then
\[
\delta(x)=[y,x]\in L^n,
\]
contradicting \(z\notin L^n\).

\subsubsection{Problem 10}

Let \(L\) be a Lie algebra, and let \(K\) be an ideal of \(L\) such that \(L/K\)
is nilpotent and such that
\[
\operatorname{ad}x|_K
\]
is nilpotent for all \(x\in L\). Prove that \(L\) is nilpotent.

\paragraph{Hint.}
Use Engel's theorem. For each \(x\in L\), the map \(\operatorname{ad}x\) preserves
\(K\). Its induced map on \(L/K\) is nilpotent because \(L/K\) is nilpotent, and
its restriction to \(K\) is nilpotent by assumption. Therefore
\(\operatorname{ad}x\) is nilpotent on all of \(L\). Engel's theorem then
implies that \(L\) is nilpotent.



\section{Semisimple Lie Algebras}

\subsection{Theorem of Lie and Cartan}

\subsubsection{Problem 1}

Let \(L=\mathfrak{sl}(V)\). Use Lie's Theorem to prove that
\(\operatorname{Rad} L=Z(L)\); conclude that \(L\) is semisimple
(cf. Exercise 2.3). Observe that \(\operatorname{Rad}L\) lies in each maximal
solvable subalgebra \(B\) of \(L\). Select a basis of \(V\) so that
\[
B=L\cap \mathfrak{t}(n,F),
\]
and notice that the transpose of \(B\) is also a maximal solvable subalgebra of
\(L\). Conclude that
\[
\operatorname{Rad}L\subset L\cap \mathfrak{d}(n,F),
\]
then that \(\operatorname{Rad}L=Z(L)\).

\paragraph{Hint.}
By Lie's Theorem, any solvable subalgebra of \(\mathfrak{gl}(V)\) can be
triangularized. Since \(\operatorname{Rad}L\) is an ideal, it is contained in
every maximal solvable subalgebra. Choose one such maximal solvable subalgebra
as upper triangular matrices, and compare it with its transpose, which is lower
triangular. Their intersection is diagonal. Finally show that any diagonal
ideal of \(\mathfrak{sl}(V)\) must lie in the center by bracketing with matrix
units \(E_{ij}\).

\subsubsection{Problem 2}

Show that the proof of Theorem 4.1 still goes through in prime characteristic,
provided \(\dim V\) is less than \(\operatorname{char} F\).

\paragraph{Hint.}
Look at where characteristic \(0\) is used in the proof of Lie's Theorem. The
key point is that certain sums of at most \(\dim V\) eigenvalues cannot vanish
for purely characteristic reasons. If \(\dim V<\operatorname{char}F\), the same
trace and eigenvalue-counting argument still works.

\subsubsection{Problem 3}

This exercise illustrates the failure of Lie's Theorem when \(F\) is allowed to
have prime characteristic \(p\). Consider the \(p\times p\) matrices
\[
x=
\begin{pmatrix}
0&1&0&\cdots&0\\
0&0&1&\cdots&0\\
\cdots&\cdots&\cdots&\cdots&\cdots\\
0&\cdots&\cdots&0&1\\
1&\cdots&\cdots&0&0
\end{pmatrix},
\qquad
y=\operatorname{diag}(0,1,2,3,\ldots,p-1).
\]
Check that \([x,y]=x\), hence that \(x\) and \(y\) span a two-dimensional
solvable subalgebra \(L\) of \(\mathfrak{gl}(p,F)\). Verify that \(x,y\) have
no common eigenvector.

\paragraph{Hint.}
Use the formula
\[
[E_{ij},y]=(j-i)E_{ij},
\]
where the indices are interpreted modulo \(p\). This gives \([x,y]=x\). To see
that there is no common eigenvector, note that \(y\) has eigenvectors the
standard basis vectors, while \(x\) cyclically permutes these basis vectors.
Thus no one-dimensional eigenspace of \(y\) is preserved by \(x\).

\subsubsection{Problem 4}

When \(p=2\), Exercise 3.3 shows that a solvable Lie algebra of endomorphisms
over a field of prime characteristic \(p\) need not have derived algebra
consisting of nilpotent endomorphisms; cf. Corollary of Theorem 4.1. For
arbitrary \(p\), construct a counterexample to Corollary C as follows. Start
with \(L\subset \mathfrak{gl}(p,F)\), as in Exercise 3. Form the vector space
direct sum
\[
M=L\oplus F^p,
\]
and make \(M\) a Lie algebra by decreeing that \(F^p\) is abelian, while \(L\)
has its usual product and acts on \(F^p\) in the given way. Verify that \(M\) is
solvable, but that its derived algebra
\[
[M,M]=Fx+F^p
\]
fails to be nilpotent.

\paragraph{Hint.}
This is a semidirect product \(M=L\ltimes F^p\). Since \(L\) is solvable and
\(F^p\) is abelian, \(M\) is solvable. Compute the derived algebra using
\[
[L,L]=Fx,\qquad [L,F^p]\subset F^p,\qquad [F^p,F^p]=0.
\]
The obstruction to nilpotency is the action of \(x\) on \(F^p\). In the example,
\(x\) is not nilpotent as an endomorphism of \(F^p\), because it cyclically
permutes the standard basis.

\subsubsection{Problem 5}

If \(x,y\in \operatorname{End}V\) commute, prove that
\[
(x+y)_s=x_s+y_s,
\qquad
(x+y)_n=x_n+y_n.
\]
Show by example that this can fail if \(x,y\) fail to commute. Show first that
\(x,y\) semisimple, respectively nilpotent, implies \(x+y\) semisimple,
respectively nilpotent.

\paragraph{Hint.}
If \(x\) and \(y\) commute, then their semisimple and nilpotent parts also
commute. For semisimple endomorphisms, use simultaneous diagonalization after
passing to an algebraic closure. For nilpotent endomorphisms, if \(x\) and \(y\)
commute, then \(x+y\) is nilpotent by the binomial theorem. For a counterexample
without commutativity, take two nilpotent \(2\times 2\) matrices whose sum is
semisimple.

\subsubsection{Problem 6}

Check formula \((*)\) at the end of (4.2).

\paragraph{Hint.}
Expand both sides using the Jordan decomposition
\[
x=x_s+x_n.
\]
Use that \(x_s\) and \(x_n\) commute. Then separate the semisimple and nilpotent
parts. The formula follows from uniqueness of the Jordan decomposition.

\subsubsection{Problem 7}

Prove the converse of Theorem 4.3.

\paragraph{Hint.}
Use Cartan's Criterion in the reverse direction. Show that if the trace
condition in Theorem 4.3 holds, then the derived algebra acts nilpotently under
the adjoint representation. Then apply Engel's Theorem to conclude solvability.

\subsubsection{Problem 8}

Note that it suffices to check the hypotheses of Theorem 4.3, or its corollary,
for \(x,y\) ranging over a basis of \(L\). For the example given in Exercise
1.2, verify solvability by using Cartan's Criterion.

\paragraph{Hint.}
The trace pairing condition is bilinear in \(x\) and \(y\), so it is enough to
check it on a basis. For the algebra in Exercise 1.2, write down the adjoint
matrices \(\operatorname{ad}x\) for basis elements \(x\), compute the relevant
traces, and verify that Cartan's Criterion applies.

\subsection{Killing form}
\subsubsection{Problem 1}

Prove that if \(L\) is nilpotent, the Killing form of \(L\) is identically
zero.

\paragraph{Hint.}
Use Engel's theorem or the fact that for nilpotent \(L\), every
\(\operatorname{ad}x\) is a nilpotent endomorphism. Then
\(\operatorname{ad}x\,\operatorname{ad}y\) has trace zero for all \(x,y\in L\).
Therefore
\[
\kappa(x,y)=\operatorname{tr}(\operatorname{ad}x\,\operatorname{ad}y)=0.
\]

\subsubsection{Problem 2}

Prove that \(L\) is solvable if and only if \([LL]\) lies in the radical of the
Killing form.

\paragraph{Hint.}
The radical of the Killing form is
\[
\operatorname{rad}\kappa
=
\{x\in L\mid \kappa(x,y)=0\text{ for all }y\in L\}.
\]
Thus \([L,L]\subset \operatorname{rad}\kappa\) means
\[
\operatorname{tr}(\operatorname{ad}x\,\operatorname{ad}y)=0
\]
for all \(x\in [L,L]\) and \(y\in L\). This is exactly Cartan's solvability
criterion.

\subsubsection{Problem 3}

Let \(L\) be the two dimensional nonabelian Lie algebra (1.4), which is
solvable. Prove that \(L\) has nontrivial Killing form.

\paragraph{Hint.}
Choose a basis \(x,y\) with
\[
[x,y]=y.
\]
Compute the adjoint matrices of \(x\) and \(y\). In particular,
\(\operatorname{ad}x\) acts nontrivially on \(Fy\). Then compute
\[
\kappa(x,x)=\operatorname{tr}((\operatorname{ad}x)^2).
\]
This trace is nonzero in characteristic \(0\), so the Killing form is not
identically zero.

\subsubsection{Problem 4}

Let \(L\) be the three dimensional solvable Lie algebra of Exercise 1.2.
Compute the radical of its Killing form.

\paragraph{Hint.}
Write down the adjoint matrices \(\operatorname{ad}x_i\) for a basis
\(x_1,x_2,x_3\) of \(L\). Then form the Killing matrix
\[
\bigl(\kappa(x_i,x_j)\bigr)_{i,j}.
\]
The radical is the nullspace of this matrix:
\[
\operatorname{rad}\kappa
=
\left\{\sum_i a_i x_i \mid
\sum_i a_i\kappa(x_i,x_j)=0
\text{ for all }j
\right\}.
\]

\subsubsection{Problem 5}

Let \(L=\mathfrak{sl}(2,F)\). Compute the basis of \(L\) dual to the standard
basis, relative to the Killing form.

\paragraph{Hint.}
Use the standard basis
\[
x=
\begin{pmatrix}
0&1\\
0&0
\end{pmatrix},
\qquad
y=
\begin{pmatrix}
0&0\\
1&0
\end{pmatrix},
\qquad
h=
\begin{pmatrix}
1&0\\
0&-1
\end{pmatrix}.
\]
Compute the values
\[
\kappa(x,y),\qquad \kappa(y,x),\qquad \kappa(h,h),
\]
and note that all other pairings among \(x,y,h\) are zero. The dual basis
\((x^\vee,y^\vee,h^\vee)\) is determined by
\[
\kappa(x,x^\vee)=1,\qquad
\kappa(y,y^\vee)=1,\qquad
\kappa(h,h^\vee)=1.
\]

\subsubsection{Problem 6}

Let \(\operatorname{char}F=p\neq 2,0\). Prove that \(L\) is semisimple if its
Killing form is nondegenerate. Show by example that the converse fails. Look at
\(\mathfrak{sl}(3,F)\) modulo its center, when \(\operatorname{char}F=3\).

\paragraph{Hint.}
If the Killing form is nondegenerate, then its radical is \(0\). Since the
radical of the Killing form is a solvable ideal, \(L\) has no nonzero solvable
ideal, hence \(L\) is semisimple. For the failure of the converse in prime
characteristic, use
\[
\kappa_{\mathfrak{sl}_n}(a,b)=2n\operatorname{tr}(ab).
\]
When \(n=3\) and \(\operatorname{char}F=3\), the scalar \(2n=6\) is zero, so the
Killing form becomes degenerate, even though
\(\mathfrak{sl}(3,F)/Z(\mathfrak{sl}(3,F))\) is semisimple.

\subsubsection{Problem 7}

Relative to the standard basis of \(\mathfrak{sl}(3,F)\), compute the
determinant of \(\kappa\).

\paragraph{Hint.}
Use the standard basis consisting of the six root vectors \(E_{ij}\), \(i\neq j\),
together with two diagonal elements, for example
\[
h_1=E_{11}-E_{22},
\qquad
h_2=E_{22}-E_{33}.
\]
Use
\[
\kappa(a,b)=6\operatorname{tr}(ab)
\]
for \(\mathfrak{sl}(3,F)\). The root-vector part pairs \(E_{ij}\) only with
\(E_{ji}\), while the diagonal part is computed from
\(\operatorname{tr}(h_i h_j)\). The determinant is the product of these two
blocks.

\subsubsection{Problem 8}

Let
\[
L=L_1\oplus\cdots\oplus L_t
\]
be the decomposition of a semisimple Lie algebra \(L\) into its simple ideals.
Show that the semisimple and nilpotent parts of \(x\in L\) are the sums of the
semisimple and nilpotent parts in the various \(L_i\) of the components of
\(x\).

\paragraph{Hint.}
Write
\[
x=x_1+\cdots+x_t,
\qquad x_i\in L_i.
\]
Since the \(L_i\) are ideals in a direct sum, one has
\[
[L_i,L_j]=0
\qquad\text{for }i\neq j.
\]
Thus the adjoint action decomposes blockwise:
\[
\operatorname{ad}x
=
\operatorname{ad}x_1\oplus\cdots\oplus \operatorname{ad}x_t.
\]
The Jordan decomposition of a block diagonal endomorphism is obtained by taking
the Jordan decomposition on each block. Use uniqueness of the abstract Jordan
decomposition in a semisimple Lie algebra.

\subsection{Complete reducibility of representation}
\subsubsection{Problem 1}

Using the standard basis for \(L=\mathfrak{sl}(2,F)\), write down the Casimir
element of the adjoint representation of \(L\). Cf. Exercise 5.5. Do the same
thing for the usual \(3\)-dimensional representation of \(\mathfrak{sl}(3,F)\),
first computing dual bases relative to the trace form.

\paragraph{Hint.}
Use the definition of the Casimir element:
\[
C_\rho=\sum_i \rho(x_i)\rho(x^i),
\]
where \(\{x_i\}\) is a basis of \(L\), and \(\{x^i\}\) is the dual basis
relative to the trace form
\[
(x,y)_\rho=\operatorname{Tr}(\rho(x)\rho(y)).
\]
For \(\mathfrak{sl}(2,F)\), use the standard basis \(x,y,h\) and compute the
dual basis from the values of the Killing form. For \(\mathfrak{sl}(3,F)\), use
the basis \(E_{ij}\), \(i\neq j\), together with
\[
h_1=E_{11}-E_{22},\qquad h_2=E_{22}-E_{33}.
\]
The dual of \(E_{ij}\) is \(E_{ji}\), while the dual diagonal basis is obtained
from the inverse of the matrix \(\bigl(\operatorname{Tr}(h_i h_j)\bigr)\).

\subsubsection{Problem 2}

Let \(V\) be an \(L\)-module. Prove that \(V\) is a direct sum of irreducible
submodules if and only if each \(L\)-submodule of \(V\) possesses a complement.

\paragraph{Hint.}
If \(V\) is completely reducible, write it as a direct sum of irreducible
submodules and choose those summands not lying in the given submodule. Conversely,
apply the complement condition repeatedly: choose an irreducible submodule
\(W\subset V\), take a complement \(V=W\oplus V_1\), then continue by induction
on \(\dim V\).

\subsubsection{Problem 3}

If \(L\) is solvable, every irreducible representation of \(L\) is
one-dimensional.

\paragraph{Hint.}
Apply Lie's Theorem to the image of \(L\) in \(\mathfrak{gl}(V)\). Since \(L\)
is solvable, there is a common eigenvector, hence a one-dimensional
\(L\)-stable subspace of \(V\). If \(V\) is irreducible, this subspace must be
all of \(V\).

\subsubsection{Problem 4}

Use Weyl's Theorem to give another proof that for \(L\) semisimple, all
\(\operatorname{Der}L\) are inner. If \(\delta\in \operatorname{Der}L\), make
the direct sum \(F\oplus L\) into an \(L\)-module via the rule
\[
x\cdot(a,y)=(0,a\delta(x)+[x,y]).
\]
Then consider a complement to the submodule \(L\).

\paragraph{Hint.}
Identify \(L\) with the submodule \(\{0\}\oplus L\subset F\oplus L\). By Weyl's
Theorem, this submodule has a one-dimensional \(L\)-stable complement. Let this
complement be generated by \((1,z)\). Since it is stable, for every \(x\in L\),
\[
x\cdot(1,z)=(0,\delta(x)+[x,z])
\]
must lie in \(F(1,z)\). Its first coordinate is \(0\), so it must be \(0\). Hence
\[
\delta(x)=-[x,z]=[z,x],
\]
so \(\delta=\operatorname{ad}z\).

\subsubsection{Problem 5}

A Lie algebra \(L\) for which \(\operatorname{Rad}L=Z(L)\) is called reductive.
Examples are \(L\) abelian, \(L\) semisimple, and \(L=\mathfrak{gl}(n,F)\).

\paragraph{(a)}
If \(L\) is reductive, then \(L\) is a completely reducible
\(\operatorname{ad}L\)-module. If \(L\neq 0\), use Weyl's Theorem. In particular,
\(L\) is the direct sum of \(Z(L)\) and \([L,L]\), with \([L,L]\) semisimple.

\paragraph{Hint.}
Since \(\operatorname{Rad}L=Z(L)\), the quotient \(L/Z(L)\) is semisimple.
The adjoint action of \(L\) factors through \(L/Z(L)\), because \(Z(L)\) acts
trivially. Apply Weyl's Theorem to the \(L/Z(L)\)-module \(L\). Then choose an
\(L\)-stable complement to \(Z(L)\). This complement is an ideal and is equal to
\([L,L]\).

\paragraph{(b)}
If \(L\) is a classical linear Lie algebra (1.2), then \(L\) is semisimple.
Cf. Exercise 1.9.

\paragraph{Hint.}
Use Exercise 1.9 to show that these classical Lie algebras are perfect:
\[
[L,L]=L.
\]
Then combine this with the earlier center computation. If a solvable radical
existed, Lie's Theorem would force it into a triangular form; using the ideal
condition and bracketing with root vectors forces it into the center. Since the
center is zero, the radical is zero.

\paragraph{(c)}
If \(L\) is a completely reducible \(\operatorname{ad}L\)-module, then \(L\) is
reductive.

\paragraph{Hint.}
Let \(R=\operatorname{Rad}L\). Since \(R\) is an ideal, it is an
\(\operatorname{ad}L\)-submodule. Complete reducibility lets one split off
\(L\)-stable complements. Use solvability of \(R\) to show that every irreducible
\(L\)-submodule contained in \(R\) is abelian, hence central. Induct on
\(\dim R\) to conclude
\[
R\subset Z(L).
\]
The reverse inclusion \(Z(L)\subset R\) is automatic, since the center is an
abelian ideal. Hence \(\operatorname{Rad}L=Z(L)\).

\paragraph{(d)}
If \(L\) is reductive, then all finite-dimensional representations of \(L\) in
which \(Z(L)\) is represented by semisimple endomorphisms are completely
reducible.

\paragraph{Hint.}
Use the decomposition
\[
L=Z(L)\oplus [L,L],
\]
where \([L,L]\) is semisimple. Since \(Z(L)\) acts by commuting semisimple
operators, decompose the representation into simultaneous eigenspaces for
\(Z(L)\). Each eigenspace is stable under \([L,L]\). Apply Weyl's Theorem to
the semisimple Lie algebra \([L,L]\) on each eigenspace.

\subsubsection{Problem 6}

Let \(L\) be a simple Lie algebra. Let \(\beta(x,y)\) and \(\gamma(x,y)\) be two
symmetric associative bilinear forms on \(L\). If \(\beta\) and \(\gamma\) are
nondegenerate, prove that \(\beta\) and \(\gamma\) are proportional. Use
Schur's Lemma.

\paragraph{Hint.}
Since \(\beta\) is nondegenerate, define a linear map \(T:L\to L\) by
\[
\gamma(x,y)=\beta(Tx,y).
\]
The associativity of \(\beta\) and \(\gamma\) implies that \(T\) commutes with
the adjoint action of \(L\). Thus \(T\) is an \(L\)-module endomorphism of the
adjoint module \(L\). Since \(L\) is simple, the adjoint module is irreducible.
By Schur's Lemma, \(T=\lambda\operatorname{id}\). Hence
\[
\gamma=\lambda\beta.
\]

\subsubsection{Problem 7}

It will be seen later on that \(\mathfrak{sl}(n,F)\) is actually simple.
Assuming this and using Exercise 6, prove that the Killing form \(\kappa\) on
\(\mathfrak{sl}(n,F)\) is related to the ordinary trace form by
\[
\kappa(x,y)=2n\operatorname{Tr}(xy).
\]

\paragraph{Hint.}
Both
\[
\kappa(x,y)
\qquad\text{and}\qquad
\operatorname{Tr}(xy)
\]
are symmetric associative nondegenerate bilinear forms on
\(\mathfrak{sl}(n,F)\). Since \(\mathfrak{sl}(n,F)\) is simple, Exercise 6
implies that
\[
\kappa(x,y)=c\operatorname{Tr}(xy)
\]
for some scalar \(c\). Determine \(c\) by evaluating both sides on a convenient
pair such as
\[
x=E_{12},\qquad y=E_{21}.
\]
A direct computation of
\(\operatorname{Tr}(\operatorname{ad}E_{12}\operatorname{ad}E_{21})\) gives
\(c=2n\).

\subsubsection{Problem 8}

If \(L\) is a Lie algebra, then \(L\) acts via \(\operatorname{ad}\) on
\((L\otimes L)^*\), which may be identified with the space of all bilinear
forms \(\beta\) on \(L\). Prove that \(\beta\) is associative if and only if
\(L\cdot \beta=0\).

\paragraph{Hint.}
The induced action on bilinear forms is
\[
(x\cdot \beta)(y,z)
=
-\beta([x,y],z)-\beta(y,[x,z]).
\]
Thus \(x\cdot\beta=0\) for all \(x\in L\) means
\[
\beta([x,y],z)+\beta(y,[x,z])=0.
\]
Since \([x,z]=-[z,x]\), this is equivalent to
\[
\beta([x,y],z)=\beta(y,[z,x]),
\]
which is precisely the associativity condition for the bilinear form.

\subsubsection{Problem 9}

Let \(L'\) be a semisimple subalgebra of a semisimple Lie algebra \(L\). If
\(x\in L'\), its Jordan decomposition in \(L'\) is also its Jordan decomposition
in \(L\).

\paragraph{Hint.}
Write the abstract Jordan decomposition of \(x\) inside \(L'\) as
\[
x=x_s+x_n.
\]
Because \(L'\) is semisimple, every finite-dimensional \(L'\)-module is
completely reducible. Apply this to the \(L'\)-module \(L\) under the adjoint
action. On every irreducible \(L'\)-submodule of \(L\), \(x_s\) acts
semisimply and \(x_n\) acts nilpotently. Hence, as endomorphisms of \(L\),
\[
\operatorname{ad}_L x
=
\operatorname{ad}_L x_s+\operatorname{ad}_L x_n
\]
is the Jordan decomposition of \(\operatorname{ad}_L x\). By uniqueness of the
abstract Jordan decomposition in the semisimple algebra \(L\), this is also the
Jordan decomposition of \(x\) in \(L\).

\subsection{Representations of \(\mathfrak{sl}(2,F)\)}

In these exercises, let \(L=\mathfrak{sl}(2,F)\).

\subsubsection{Problem 1}

Use Lie's Theorem to prove the existence of a maximal vector in an arbitrary
finite-dimensional \(L\)-module. Look at the subalgebra \(B\) spanned by
\(h\) and \(x\).

\paragraph{Hint.}
The subalgebra \(B=Fh+Fx\) is solvable since \([h,x]=2x\). By Lie's Theorem,
there is a common eigenvector for \(B\). In particular, it is an eigenvector
for \(h\), and \(x\) acts by a scalar. Use the relation \([h,x]=2x\) to show
that this scalar must be \(0\). Hence this vector is killed by \(x\), so it is
a maximal vector.

\subsubsection{Problem 2}

Let \(M=\mathfrak{sl}(3,F)\) contain a copy of \(L\) in its upper left-hand
\(2\times 2\) position. Write \(M\) as a direct sum of irreducible
\(L\)-submodules, where \(M\) is viewed as an \(L\)-module via the adjoint
representation:
\[
V(0)\oplus V(1)\oplus V(1)\oplus V(2).
\]

\paragraph{Hint.}
Decompose the matrix units \(E_{ij}\) according to how the upper-left
\(\mathfrak{sl}(2,F)\) acts by commutator. The upper-left \(2\times 2\)
trace-zero block gives the adjoint module \(V(2)\). The scalar diagonal part
commuting with this copy of \(L\) gives \(V(0)\). The first two entries of the
third column and the first two entries of the third row each form a copy of
the standard \(2\)-dimensional module \(V(1)\).

\subsubsection{Problem 3}

Verify that formulas (a)--(c) of Lemma 7.2 do define an irreducible
representation of \(L\). To show that they define a representation, it suffices
to check that the matrices corresponding to \(x,y,h\) satisfy the same
structural equations as \(x,y,h\).

\paragraph{Hint.}
Check the three defining relations of \(\mathfrak{sl}(2,F)\):
\[
[h,x]=2x,\qquad [h,y]=-2y,\qquad [x,y]=h.
\]
Apply both sides to each basis vector \(v_i\). For irreducibility, take a
nonzero submodule and choose a vector in it with maximal or minimal occurring
index. Use the action of \(x\) and \(y\) to force the submodule to contain a
maximal vector, hence \(v_0\), and then all basis vectors.

\subsubsection{Problem 4}

The irreducible representation of \(L\) of highest weight \(m\) can also be
realized naturally as follows. Let \(X,Y\) be a basis for the two-dimensional
vector space \(F^2\), on which \(L\) acts as usual. Let
\[
\mathcal{R}=F[X,Y]
\]
be the polynomial algebra in two variables, and extend the action of \(L\) to
\(\mathcal{R}\) by the derivation rule
\[
z\cdot fg=(z\cdot f)g+f(z\cdot g),
\qquad z\in L,\ f,g\in\mathcal{R}.
\]
Show that this extension is well-defined and that \(\mathcal{R}\) becomes an
\(L\)-module. Then show that the subspace of homogeneous polynomials of degree
\(m\), with basis
\[
X^m,\ X^{m-1}Y,\ \ldots,\ XY^{m-1},\ Y^m,
\]
is invariant under \(L\) and irreducible.

\paragraph{Hint.}
First check the action on monomials using the Leibniz rule. The operators
coming from \(x,y,h\) are derivations of the polynomial algebra and satisfy
the same \(\mathfrak{sl}(2,F)\)-relations. Since \(x,y,h\) act by differential
operators preserving total degree, the homogeneous degree \(m\) part is
stable. Identify \(X^m\) as a maximal vector of weight \(m\), and use repeated
application of \(y\) to generate the whole degree \(m\) subspace.

\subsubsection{Problem 5}

Suppose \(\operatorname{char}F=p>0\) and \(L=\mathfrak{sl}(2,F)\). Prove that
the representation \(V(m)\) of \(L\) constructed as in Exercise 3 or 4 is
irreducible so long as the highest weight \(m\) is strictly less than \(p\),
but reducible when \(m=p\).

\paragraph{Hint.}
In the usual basis \(v_0,\ldots,v_m\), the coefficients in the action of \(x\)
and \(y\) involve integers such as \(i\) and \(m-i+1\). If \(m<p\), none of the
necessary coefficients vanish modulo \(p\), so the standard irreducibility
argument still works. When \(m=p\), one of these coefficients becomes \(0\) in
\(F\), producing a proper nonzero submodule.

\subsubsection{Problem 6}

Decompose the tensor product of the two \(L\)-modules \(V(3)\), \(V(7)\) into
the sum of irreducible submodules:
\[
V(4)\oplus V(6)\oplus V(8)\oplus V(10).
\]
Try to develop a general formula for the decomposition of \(V(m)\otimes V(n)\).

\paragraph{Hint.}
Use highest weights. The largest possible highest weight in
\(V(m)\otimes V(n)\) is \(m+n\). After removing \(V(m+n)\), look for the next
maximal vector. The general Clebsch--Gordan decomposition is
\[
V(m)\otimes V(n)
\cong
V(m+n)\oplus V(m+n-2)\oplus\cdots\oplus V(|m-n|).
\]
For \(m=3\) and \(n=7\), this gives
\[
V(10)\oplus V(8)\oplus V(6)\oplus V(4).
\]

\subsubsection{Problem 7}

In this exercise we construct certain infinite-dimensional \(L\)-modules. Let
\(\lambda\in F\) be an arbitrary scalar. Let \(Z(\lambda)\) be a vector space
over \(F\) with countably infinite basis
\[
v_0,v_1,v_2,\ldots .
\]

\paragraph{(a)}
Prove that formulas (a)--(c) of Lemma 7.2 define an \(L\)-module structure on
\(Z(\lambda)\), and that every nonzero \(L\)-submodule of \(Z(\lambda)\)
contains at least one maximal vector.

\paragraph{Hint.}
As in the finite-dimensional case, verify the defining relations
\[
[h,x]=2x,\qquad [h,y]=-2y,\qquad [x,y]=h
\]
on every basis vector \(v_i\). For the maximal-vector statement, take a
nonzero vector in a submodule with minimal occurring index. Applying \(x\)
lowers the index, so minimality forces the resulting vector to be killed by
\(x\).

\paragraph{(b)}
Suppose \(\lambda+1=i\) is a nonnegative integer. Prove that \(v_i\) is a
maximal vector, for example \(\lambda=-1\), \(i=0\). This induces an
\(L\)-module homomorphism
\[
\phi:Z(\mu)\longrightarrow Z(\lambda),
\qquad
\mu=\lambda-2i,
\]
sending \(v_0\) to \(v_i\). Show that \(\phi\) is a monomorphism, and that
\(\operatorname{Im}\phi\) and \(Z(\lambda)/\operatorname{Im}\phi\) are both
irreducible \(L\)-modules, but \(Z(\lambda)\) fails to be completely reducible
when \(i>0\).

\paragraph{Hint.}
Use the formula for \(xv_j\). The coefficient of \(v_{j-1}\) vanishes exactly
when \(j=\lambda+1\), so \(v_i\) is killed by \(x\). The submodule generated by
\(v_i\) has highest weight \(\lambda-2i\), giving the map from \(Z(\mu)\).
Injectivity follows because the images of the basis vectors remain linearly
independent. The quotient has basis represented by \(v_0,\ldots,v_{i-1}\) and
is the finite-dimensional irreducible highest-weight module.

\paragraph{(c)}
Suppose \(\lambda+1\) is not a nonnegative integer. Prove that \(Z(\lambda)\)
is irreducible.

\paragraph{Hint.}
By part (a), any nonzero submodule contains a maximal vector. Under the present
assumption, the only possible maximal vector is \(v_0\). Hence any nonzero
submodule contains \(v_0\), and repeated application of \(y\) generates all
basis vectors \(v_0,v_1,v_2,\ldots\). Therefore the submodule is all of
\(Z(\lambda)\).

\subsection{Root Space Decomposition}

\subsubsection{Problem 1}

If \(L\) is a classical linear Lie algebra of type \(A_\ell,B_\ell,C_\ell\), or
\(D_\ell\) (see (1.2)), prove that the set of all diagonal matrices in \(L\) is
a maximal toral subalgebra, of dimension \(\ell\). Cf. Exercise 2.8.

\paragraph{Hint.}
First write down the diagonal subalgebra \(H\) explicitly. For type \(A_\ell\),
it is the trace-zero diagonal subalgebra of \(\mathfrak{sl}(\ell+1,F)\). For
types \(B_\ell,C_\ell,D_\ell\), it consists of diagonal matrices of the form
\[
\operatorname{diag}(a_1,\ldots,a_\ell,0,-a_1,\ldots,-a_\ell)
\]
or
\[
\operatorname{diag}(a_1,\ldots,a_\ell,-a_1,\ldots,-a_\ell),
\]
depending on the type. These matrices commute and act semisimply by
\(\operatorname{ad}\). To prove maximality, show that the centralizer of \(H\)
inside \(L\) is exactly \(H\). Hence no larger toral subalgebra can contain
\(H\).

\subsubsection{Problem 2}

For each algebra in Exercise 1, determine the roots and root spaces. How are
the various \(h_\alpha\) expressed in terms of the basis for \(H\) given in
(1.2)?

\paragraph{Hint.}
Let \(\varepsilon_i\in H^*\) be the functional extracting the \(i\)-th diagonal
coordinate. The roots are
\[
A_\ell:\quad \varepsilon_i-\varepsilon_j,\quad i\neq j,
\]
\[
B_\ell:\quad \pm \varepsilon_i\pm \varepsilon_j,\ i\neq j,\quad
\pm\varepsilon_i,
\]
\[
C_\ell:\quad \pm \varepsilon_i\pm \varepsilon_j,\ i\neq j,\quad
\pm 2\varepsilon_i,
\]
\[
D_\ell:\quad \pm \varepsilon_i\pm \varepsilon_j,\quad i\neq j.
\]
For type \(A_\ell\), the root space for \(\varepsilon_i-\varepsilon_j\) is
\(F E_{ij}\). For the other classical types, the root vectors are suitable
linear combinations of matrix units forced by the defining bilinear form. The
element \(h_\alpha\) is obtained from the condition
\[
\beta(h_\alpha,h)=\alpha(h)
\qquad \text{for all } h\in H,
\]
where \(\beta\) is the Killing form, or equivalently from the standard
\(\mathfrak{sl}_2\)-triple \([x_\alpha,x_{-\alpha}]=h_\alpha\).

\subsubsection{Problem 3}

If \(L\) is of classical type, compute explicitly the restriction of the
Killing form to the maximal toral subalgebra described in Exercise 1.

\paragraph{Hint.}
Use the trace formulas
\[
\kappa_{\mathfrak{sl}_n}(x,y)=2n\operatorname{Tr}(xy),
\qquad
\kappa_{\mathfrak{so}_n}(x,y)=(n-2)\operatorname{Tr}(xy),
\]
and
\[
\kappa_{\mathfrak{sp}_{2\ell}}(x,y)=2(\ell+1)\operatorname{Tr}(xy).
\]
Then substitute diagonal elements
\[
h=\operatorname{diag}(a_1,\ldots),\qquad
h'=\operatorname{diag}(b_1,\ldots).
\]
For example, in type \(A_\ell\), with \(n=\ell+1\),
\[
\kappa(h,h')=2n\sum_{i=1}^n a_i b_i,
\qquad \sum_i a_i=\sum_i b_i=0.
\]
For types \(B,C,D\), the trace of \(hh'\) is twice \(\sum_i a_i b_i\).

\subsubsection{Problem 4}

If \(L=\mathfrak{sl}(2,F)\), prove that each maximal toral subalgebra is one
dimensional.

\paragraph{Hint.}
A toral subalgebra consists of semisimple elements and is abelian. If
\(0\neq h\in H\) is semisimple, then after choosing a suitable basis it is
conjugate to a scalar multiple of
\[
\begin{pmatrix}
1&0\\
0&-1
\end{pmatrix}.
\]
Its centralizer in \(\mathfrak{sl}(2,F)\) is one-dimensional. Since every
element of a toral subalgebra containing \(h\) must commute with \(h\), the
whole toral subalgebra lies in \(C_L(h)\). Hence it has dimension \(1\).

\subsubsection{Problem 5}

If \(L\) is semisimple and \(H\) is a maximal toral subalgebra, prove that \(H\)
is self-normalizing, i.e.
\[
H=N_L(H).
\]

\paragraph{Hint.}
Use the root space decomposition
\[
L=H\oplus \bigoplus_{\alpha\in \Phi} L_\alpha.
\]
Let \(x\in N_L(H)\), and write
\[
x=h+\sum_{\alpha\in\Phi}x_\alpha,
\qquad h\in H,\ x_\alpha\in L_\alpha.
\]
For every \(t\in H\), the condition \(x\in N_L(H)\) gives
\[
[t,x]\in H.
\]
But
\[
[t,x]=\sum_{\alpha\in\Phi}\alpha(t)x_\alpha.
\]
Since the summands lie in root spaces and not in \(H\), all \(x_\alpha\) must
be zero. Therefore \(x\in H\), so \(N_L(H)=H\).

\subsubsection{Problem 6}

Compute the basis of \(\mathfrak{sl}(n,F)\) which is dual, via the Killing
form, to the standard basis. Cf. Exercise 5.5.

\paragraph{Hint.}
Use the standard basis
\[
E_{ij}\quad (i\neq j),\qquad
h_i=E_{ii}-E_{i+1,i+1},\quad 1\leq i\leq n-1.
\]
Since
\[
\kappa(x,y)=2n\operatorname{Tr}(xy),
\]
the element dual to \(E_{ij}\) is
\[
\frac{1}{2n}E_{ji}.
\]
For the diagonal part, compute the matrix
\[
\bigl(\kappa(h_i,h_j)\bigr)_{i,j}
=
2n\bigl(\operatorname{Tr}(h_i h_j)\bigr)_{i,j}.
\]
The trace matrix is the Cartan matrix of type \(A_{n-1}\). Therefore the dual
basis to the \(h_i\)'s is obtained by multiplying by the inverse Cartan matrix.

\subsubsection{Problem 7}

Let \(L\) be semisimple and \(H\) a maximal toral subalgebra. If \(h\in H\),
prove that \(C_L(h)\) is reductive in the sense of Exercise 6.5. Prove that
\(H\) contains elements \(h\) for which \(C_L(h)=H\). For which \(h\in
\mathfrak{sl}(n,F)\) is this true?

\paragraph{Hint.}
Using the root space decomposition, one has
\[
C_L(h)
=
H\oplus \bigoplus_{\alpha(h)=0} L_\alpha.
\]
This is a Lie subalgebra. Its center contains the part of \(H\) killed by the
roots appearing in the centralizer, and its derived algebra is generated by
the root spaces with \(\alpha(h)=0\). This gives a reductive decomposition.

For \(C_L(h)=H\), choose \(h\in H\) such that
\[
\alpha(h)\neq 0
\qquad \text{for all } \alpha\in\Phi.
\]
Such an \(h\) is called regular. In \(\mathfrak{sl}(n,F)\), this means that
\(h=\operatorname{diag}(a_1,\ldots,a_n)\), with \(\sum_i a_i=0\), has pairwise
distinct diagonal entries:
\[
a_i\neq a_j
\qquad \text{for } i\neq j.
\]

\subsubsection{Problem 8}

For \(\mathfrak{sl}(n,F)\) and other classical algebras, calculate explicitly
the root strings and Cartan integers. In particular, prove that all Cartan
integers
\[
\frac{2(\alpha,\beta)}{(\beta,\beta)},
\qquad \alpha\neq \pm\beta,
\]
for \(\mathfrak{sl}(n,F)\) are \(0,\pm 1\).

\paragraph{Hint.}
For \(\mathfrak{sl}(n,F)\), write every root as
\[
\alpha=\varepsilon_i-\varepsilon_j.
\]
Use the standard inner product on the hyperplane
\[
\sum_i \varepsilon_i=0.
\]
Then compute
\[
(\varepsilon_i-\varepsilon_j,\varepsilon_k-\varepsilon_l)
=
\delta_{ik}-\delta_{il}-\delta_{jk}+\delta_{jl}.
\]
Since every root has squared length \(2\), the Cartan integer is just this
inner product. For \(\alpha\neq \pm \beta\), the only possible values are
\(0,\pm 1\). The root string through \(\alpha\) in the \(\beta\)-direction is
found by checking which of
\[
\alpha-r\beta,\ldots,\alpha,\ldots,\alpha+q\beta
\]
are again roots.

\subsubsection{Problem 9}

Prove that every three dimensional semisimple Lie algebra has the same root
system as \(\mathfrak{sl}(2,F)\), hence is isomorphic to \(\mathfrak{sl}(2,F)\).

\paragraph{Hint.}
Let \(H\) be a maximal toral subalgebra. Since \(L\) is semisimple,
\[
L=H\oplus \bigoplus_{\alpha\in\Phi}L_\alpha.
\]
The roots occur in pairs \(\alpha,-\alpha\), and each root space has dimension
\(1\). Since \(\dim L=3\), the only possibility is
\[
\dim H=1,
\qquad
\Phi=\{\alpha,-\alpha\}.
\]
This is the root system of type \(A_1\). Choose root vectors
\(x\in L_\alpha\), \(y\in L_{-\alpha}\), and \(h=[x,y]\in H\). After rescaling,
the relations become
\[
[h,x]=2x,\qquad [h,y]=-2y,\qquad [x,y]=h,
\]
so \(L\cong \mathfrak{sl}(2,F)\).

\subsubsection{Problem 10}

Prove that no four, five, or seven dimensional semisimple Lie algebras exist.

\paragraph{Hint.}
A semisimple Lie algebra is a direct sum of simple ideals. The smallest
nonzero simple Lie algebra is \(\mathfrak{sl}(2,F)\), of dimension \(3\). If a
simple Lie algebra has rank \(1\), it is \(\mathfrak{sl}(2,F)\). If its rank is
at least \(2\), its root system has at least six roots, so its dimension is at
least
\[
2+6=8.
\]
Therefore the possible small dimensions of semisimple Lie algebras begin with
\[
3,\qquad 3+3=6,\qquad 8,\ldots
\]
Thus dimensions \(4,5,7\) cannot occur.

\subsubsection{Problem 11}

If \((\alpha,\beta)>0\), prove that
\[
\alpha-\beta\in \Phi
\qquad (\alpha,\beta\in\Phi).
\]
Is the converse true?

\paragraph{Hint.}
Assume \(\alpha\neq \beta\). Consider the \(\beta\)-string through \(\alpha\):
\[
\alpha-r\beta,\ldots,\alpha,\ldots,\alpha+q\beta.
\]
The root string theorem gives
\[
r-q=\frac{2(\alpha,\beta)}{(\beta,\beta)}.
\]
If \((\alpha,\beta)>0\), then \(r-q>0\), hence \(r>0\). Therefore
\(\alpha-\beta\in \Phi\).

The converse is false. In type \(B_2\), take
\[
\alpha=\varepsilon_1,\qquad \beta=\varepsilon_2.
\]
Then
\[
\alpha-\beta=\varepsilon_1-\varepsilon_2
\]
is a root, but
\[
(\alpha,\beta)=0.
\]

\section{Root Systems}

\subsection{Axioms}

Unless otherwise specified, \(\Phi\) denotes a root system in \(E\), with Weyl
group \(\mathcal W\).

\subsubsection{Problem 1}

Let \(E'\) be a subspace of \(E\). If a reflection \(\sigma_\alpha\) leaves
\(E'\) invariant, prove that either \(\alpha\in E'\) or \(E'\subset P_\alpha\).

\paragraph{Hint.}
Use the formula
\[
\sigma_\alpha(v)=v-\frac{2(v,\alpha)}{(\alpha,\alpha)}\alpha .
\]
If some \(v\in E'\) satisfies \((v,\alpha)\neq 0\), then
\(v-\sigma_\alpha(v)\) is a nonzero scalar multiple of \(\alpha\), so
\(\alpha\in E'\). Otherwise \((v,\alpha)=0\) for all \(v\in E'\), hence
\(E'\subset P_\alpha\).

\subsubsection{Problem 2}

Prove that \(\Phi^\vee\) is a root system in \(E\), whose Weyl group is
naturally isomorphic to \(\mathcal W\). Show also that
\[
\langle \alpha^\vee,\beta^\vee\rangle
=
\langle \beta,\alpha\rangle,
\]
and draw a picture of \(\Phi^\vee\) in the cases \(A_1,A_2,B_2,G_2\).

\paragraph{Hint.}
Recall that
\[
\alpha^\vee=\frac{2\alpha}{(\alpha,\alpha)}.
\]
The reflection hyperplane of \(\alpha^\vee\) is the same as that of \(\alpha\),
and
\[
\sigma_{\alpha^\vee}=\sigma_\alpha.
\]
Thus the Weyl group does not change. The axioms for \(\Phi^\vee\) follow by
checking that reflections preserve coroots and that the possible scalar
multiples are still only \(\pm\alpha^\vee\). For the Cartan integer, compute
directly:
\[
\langle \alpha^\vee,\beta^\vee\rangle
=
\frac{2(\alpha^\vee,\beta^\vee)}{(\beta^\vee,\beta^\vee)}
=
\frac{2(\beta,\alpha)}{(\alpha,\alpha)}
=
\langle \beta,\alpha\rangle.
\]
Geometrically, passing to the dual root system interchanges long and short
roots. Thus \(A_1\) and \(A_2\) are unchanged, \(B_2^\vee\) is \(C_2\), and
\(G_2^\vee\) is again \(G_2\) with long and short roots interchanged.

\subsubsection{Problem 3}

In Table 1, show that the order of \(\sigma_\alpha\sigma_\beta\) in
\(\mathcal W\) is respectively
\[
2,\ 3,\ 4,\ 6
\]
when
\[
\theta=\pi/2,\ \pi/3\text{ or }2\pi/3,\ \pi/4\text{ or }3\pi/4,\ 
\pi/6\text{ or }5\pi/6.
\]
Note that \(\sigma_\alpha\sigma_\beta\) is a rotation through \(2\theta\).

\paragraph{Hint.}
In the plane spanned by \(\alpha\) and \(\beta\), the product of two reflections
across lines with angle \(\theta\) is a rotation through angle \(2\theta\). The
order of this rotation is the smallest positive integer \(m\) such that
\[
m(2\theta)\in 2\pi\mathbb Z.
\]
Thus
\[
\theta=\pi/2,\pi/3,\pi/4,\pi/6
\]
give orders \(2,3,4,6\), respectively. The obtuse-angle cases give the same
orders because replacing \(\theta\) by \(\pi-\theta\) changes the rotation
angle by sign modulo \(2\pi\).

\subsubsection{Problem 4}

Prove that the respective Weyl groups of
\[
A_1\times A_1,\quad A_2,\quad B_2,\quad G_2
\]
are dihedral of orders
\[
4,\ 6,\ 8,\ 12.
\]
If \(\Phi\) is any root system of rank \(2\), prove that its Weyl group must be
one of these.

\paragraph{Hint.}
The Weyl group is generated by the two simple reflections
\(\sigma_\alpha,\sigma_\beta\). From Problem 3, the product
\(\sigma_\alpha\sigma_\beta\) has order \(m=2,3,4,6\), depending on the angle
between the simple roots. Hence
\[
\mathcal W=\langle \sigma_\alpha,\sigma_\beta
\mid \sigma_\alpha^2=\sigma_\beta^2=1,\ 
(\sigma_\alpha\sigma_\beta)^m=1\rangle
\]
is a dihedral group of order \(2m\). The root-system axioms restrict rank-two
angles to exactly these possibilities, so the only possible rank-two Weyl
groups are the four listed above.

\subsubsection{Problem 5}

Show by example that \(\alpha-\beta\) may be a root even when
\[
(\alpha,\beta)\leq 0.
\]
Cf. Lemma 9.4.

\paragraph{Hint.}
Use type \(B_2\). Let
\[
\alpha=\varepsilon_1,\qquad \beta=\varepsilon_2.
\]
Then
\[
(\alpha,\beta)=0,
\]
but
\[
\alpha-\beta=\varepsilon_1-\varepsilon_2
\]
is a root of \(B_2\). This shows that the implication
\[
\alpha-\beta\in\Phi \Longrightarrow (\alpha,\beta)>0
\]
is false.

\subsubsection{Problem 6}

Prove that \(\mathcal W\) is a normal subgroup of
\[
\operatorname{Aut}\Phi,
\]
where \(\operatorname{Aut}\Phi\) denotes the group of all isomorphisms of
\(\Phi\) onto itself.

\paragraph{Hint.}
If \(g\in \operatorname{Aut}\Phi\), then \(g\) sends roots to roots and
preserves the reflection structure. Show that
\[
g\sigma_\alpha g^{-1}=\sigma_{g(\alpha)}.
\]
Since \(\mathcal W\) is generated by the reflections \(\sigma_\alpha\), this
implies
\[
g\mathcal W g^{-1}\subset \mathcal W.
\]
Applying the same argument to \(g^{-1}\) gives equality, hence
\(\mathcal W\triangleleft \operatorname{Aut}\Phi\).

\subsubsection{Problem 7}

Let \(\alpha,\beta\in\Phi\) span a subspace \(E'\) of \(E\). Prove that
\(E'\cap\Phi\) is a root system in \(E'\). Prove similarly that
\[
\Phi\cap(\mathbb Z\alpha+\mathbb Z\beta)
\]
is a root system in \(E'\). Must this coincide with \(E'\cap\Phi\)?

More generally, let \(\Phi'\) be a nonempty subset of \(\Phi\) such that
\[
\Phi'=-\Phi',
\]
and such that whenever \(\alpha,\beta\in\Phi'\) and
\(\alpha+\beta\in\Phi\), one has
\[
\alpha+\beta\in\Phi'.
\]
Prove that \(\Phi'\) is a root system in the subspace of \(E\) it spans. Use
Table 1.

\paragraph{Hint.}
For \(E'\cap\Phi\), check the axioms inside the plane \(E'\). If
\(\gamma\in E'\cap\Phi\), then \(\sigma_\gamma\) preserves \(E'\), so it
preserves \(E'\cap\Phi\). The integrality and scalar axioms are inherited from
\(\Phi\).

For \(\Phi\cap(\mathbb Z\alpha+\mathbb Z\beta)\), use that reflections send
roots to integral linear combinations of roots:
\[
\sigma_\gamma(\delta)=\delta-\langle \delta,\gamma\rangle\gamma.
\]
Thus the lattice \(\mathbb Z\alpha+\mathbb Z\beta\) is preserved by the
relevant reflections.

For the general statement, the condition on sums says that \(\Phi'\) is closed
under root strings. In each rank-two subsystem from Table 1, this closure under
\(\alpha+\beta\) and symmetry under negatives is enough to ensure stability
under the reflections. Hence \(\Phi'\) satisfies the root-system axioms in its
span.

\subsubsection{Problem 8}

Compute root strings in \(G_2\) to verify the relation
\[
r-q=\langle \beta,\alpha\rangle.
\]

\paragraph{Hint.}
Choose simple roots \(\alpha,\beta\) for \(G_2\), with \(\alpha\) short and
\(\beta\) long. The positive roots are
\[
\alpha,\quad \beta,\quad \alpha+\beta,\quad 2\alpha+\beta,\quad
3\alpha+\beta,\quad 3\alpha+2\beta.
\]
For each pair of roots, list the roots in the \(\alpha\)-string through
\(\beta\):
\[
\beta-r\alpha,\ldots,\beta,\ldots,\beta+q\alpha.
\]
For example, the \(\alpha\)-string through \(\beta\) is
\[
\beta,\quad \beta+\alpha,\quad \beta+2\alpha,\quad \beta+3\alpha,
\]
so \(r=0\), \(q=3\), and therefore
\[
r-q=-3=\langle \beta,\alpha\rangle.
\]
Check the other strings similarly.

\subsubsection{Problem 9}

Let \(\Phi\) be a set of vectors in a Euclidean space \(E\), satisfying only
\((R1)\), \((R3)\), and \((R4)\). Prove that the only possible multiples of
\(\alpha\in\Phi\) which can be in \(\Phi\) are
\[
\pm\frac{1}{2}\alpha,\qquad \pm\alpha,\qquad \pm2\alpha.
\]
Verify that
\[
\{\alpha\in\Phi\mid 2\alpha\notin\Phi\}
\]
is a root system.

\paragraph{Hint.}
Suppose \(c\alpha\in\Phi\). Use the integrality axiom in both directions:
\[
\langle c\alpha,\alpha\rangle=2c\in\mathbb Z,
\qquad
\langle \alpha,c\alpha\rangle=\frac{2}{c}\in\mathbb Z.
\]
Thus both \(2c\) and \(2/c\) are integers, forcing
\[
c\in\left\{\pm\frac12,\pm1,\pm2\right\}.
\]

For the second claim, remove precisely the roots which are twice another root.
The remaining set is reduced by construction. Check that it is still stable
under reflections: if \(2\sigma_\alpha(\beta)\) were a root, then applying
\(\sigma_\alpha^{-1}\) would imply that \(2\beta\) is a root. Hence the subset
satisfies the reduced root-system axioms.


\subsection{Simple Roots and Weyl Groups}

Unless otherwise specified, \(\Phi\) denotes a root system in \(E\), with Weyl
group \(\mathcal W\).

\subsubsection{Problem 1}

Let \(\Phi^\vee\) be the dual system of \(\Phi\), and let
\[
\Delta^\vee=\{\alpha^\vee\mid \alpha\in \Delta\}.
\]
Prove that \(\Delta^\vee\) is a base of \(\Phi^\vee\). Compare Weyl chambers of
\(\Phi\) and \(\Phi^\vee\).

\paragraph{Hint.}
Use that
\[
P_{\alpha^\vee}=P_\alpha
\]
because \(\alpha^\vee\) is a positive scalar multiple of \(\alpha\). Hence
\(\Phi\) and \(\Phi^\vee\) have the same reflecting hyperplanes and the same
Weyl chambers. The chamber defined by
\[
(\lambda,\alpha)>0\qquad(\alpha\in\Delta)
\]
is also defined by
\[
(\lambda,\alpha^\vee)>0\qquad(\alpha^\vee\in\Delta^\vee).
\]
Therefore the walls of the same fundamental chamber are indexed by
\(\Delta^\vee\), so \(\Delta^\vee\) is a base of \(\Phi^\vee\).

\subsubsection{Problem 2}

If \(\Delta\) is a base of \(\Phi\), prove that the set
\[
(\mathbb Z\alpha+\mathbb Z\beta)\cap\Phi
\]
where \(\alpha\neq\beta\) are in \(\Delta\), is a root system of rank \(2\) in
the subspace of \(E\) spanned by \(\alpha,\beta\). Cf. Exercise 9.7. Generalize
to an arbitrary subset of \(\Delta\).

\paragraph{Hint.}
Let
\[
E_{\alpha\beta}=\mathbb R\alpha+\mathbb R\beta.
\]
The set
\[
\Phi_{\alpha\beta}:=(\mathbb Z\alpha+\mathbb Z\beta)\cap\Phi
\]
is finite, spans \(E_{\alpha\beta}\), and inherits the integrality condition
from \(\Phi\). For \(\gamma,\delta\in\Phi_{\alpha\beta}\), use
\[
\sigma_\gamma(\delta)=\delta-\langle\delta,\gamma\rangle\gamma.
\]
Since \(\langle\delta,\gamma\rangle\in\mathbb Z\), the reflection
\(\sigma_\gamma\) preserves \(\mathbb Z\alpha+\mathbb Z\beta\). Hence it
preserves \(\Phi_{\alpha\beta}\). The same argument works for
\[
\Phi_J:=\Phi\cap\sum_{\alpha\in J}\mathbb Z\alpha
\]
where \(J\subset\Delta\).

\subsubsection{Problem 3}

Prove that each root system of rank \(2\) is isomorphic to one of those listed
in (9.3).

\paragraph{Hint.}
Choose two simple roots \(\alpha,\beta\). The root-system axioms imply
\[
\langle\alpha,\beta\rangle\langle\beta,\alpha\rangle
=
4\cos^2\theta,
\]
where \(\theta\) is the angle between \(\alpha\) and \(\beta\). Since both
Cartan integers are nonpositive integers for simple roots, the only possible
products are
\[
0,\ 1,\ 2,\ 3.
\]
These give respectively
\[
A_1\times A_1,\qquad A_2,\qquad B_2,\qquad G_2.
\]
Then use the root-string property to recover all roots in each case.

\subsubsection{Problem 4}

Verify the Corollary of Lemma 10.2A directly for \(G_2\).

\paragraph{Hint.}
Take simple roots \(\alpha,\beta\), where \(\alpha\) is short and \(\beta\) is
long. The positive roots are
\[
\alpha,\quad \beta,\quad \alpha+\beta,\quad 2\alpha+\beta,\quad
3\alpha+\beta,\quad 3\alpha+2\beta.
\]
Check directly that every positive root is a nonnegative integral combination
of the simple roots. For each positive root which is not simple, verify that
subtracting a suitable simple root still gives a root. This is the rank-two
verification of the corollary.

\subsubsection{Problem 5}

If \(\sigma\in\mathcal W\) can be written as a product of \(t\) simple
reflections, prove that \(t\) has the same parity as \(\ell(\sigma)\).

\paragraph{Hint.}
Use the fact that multiplication by a simple reflection changes length by
exactly \(1\):
\[
\ell(\sigma\tau_\alpha)=\ell(\sigma)\pm 1.
\]
Starting from the identity, whose length is \(0\), after multiplying by \(t\)
simple reflections the length has the same parity as \(t\). Therefore
\[
t\equiv \ell(\sigma)\pmod 2.
\]

\subsubsection{Problem 6}

Define a function
\[
\operatorname{sn}:\mathcal W\to\{\pm 1\}
\]
by
\[
\operatorname{sn}(\sigma)=(-1)^{\ell(\sigma)}.
\]
Prove that \(\operatorname{sn}\) is a homomorphism. Cf. the case \(A_2\), where
\(\mathcal W\) is isomorphic to the symmetric group \(S_3\).

\paragraph{Hint.}
Write reduced expressions
\[
\sigma=s_1\cdots s_r,\qquad \tau=t_1\cdots t_m.
\]
Then \(\sigma\tau\) is expressed as a product of \(r+m\) simple reflections.
By Problem 5,
\[
\ell(\sigma\tau)\equiv r+m=\ell(\sigma)+\ell(\tau)\pmod 2.
\]
Hence
\[
\operatorname{sn}(\sigma\tau)
=
(-1)^{\ell(\sigma\tau)}
=
(-1)^{\ell(\sigma)}(-1)^{\ell(\tau)}.
\]

\subsubsection{Problem 7}

Prove that the intersection of the ``positive'' open half-spaces associated
with any basis \(\gamma_1,\ldots,\gamma_\ell\) of \(E\) is nonvoid. If
\(\delta_i\) is the projection of \(\gamma_i\) on the orthogonal complement of
the subspace spanned by all basis vectors except \(\gamma_i\), consider
\[
\gamma=\sum r_i\delta_i
\]
when all \(r_i>0\).

\paragraph{Hint.}
The vectors \(\delta_i\) satisfy
\[
(\delta_i,\gamma_j)=0\quad(j\neq i),
\qquad
(\delta_i,\gamma_i)>0.
\]
Therefore, if
\[
\gamma=\sum_i r_i\delta_i
\]
with all \(r_i>0\), then
\[
(\gamma,\gamma_j)=r_j(\delta_j,\gamma_j)>0
\]
for every \(j\). Hence \(\gamma\) lies in all the required positive open
half-spaces.

\subsubsection{Problem 8}

Let \(\Delta\) be a base of \(\Phi\), and let \(\alpha\neq\beta\) be simple
roots. Let \(\Phi_{\alpha\beta}\) be the rank two root system in
\[
E_{\alpha\beta}=\mathbb R\alpha+\mathbb R\beta.
\]
The Weyl group \(\mathcal W_{\alpha\beta}\) of \(\Phi_{\alpha\beta}\) is
generated by the restrictions \(\tau_\alpha,\tau_\beta\) to \(E_{\alpha\beta}\)
of \(\sigma_\alpha,\sigma_\beta\), and \(\mathcal W_{\alpha\beta}\) may be
viewed as a subgroup of \(\mathcal W\). Prove that the length of an element of
\(\mathcal W_{\alpha\beta}\), relative to \(\tau_\alpha,\tau_\beta\), coincides
with the length of the corresponding element of \(\mathcal W\).

\paragraph{Hint.}
Use the inversion-set formula
\[
\ell(w)=\#\{\gamma\in\Phi^+\mid w\gamma\in\Phi^-\}.
\]
For \(w\in\mathcal W_{\alpha\beta}\), the roots whose signs can change are
exactly the roots in
\[
\Phi_{\alpha\beta}=(\mathbb Z\alpha+\mathbb Z\beta)\cap\Phi.
\]
Thus the inversion set computed in the full root system is the same as the
inversion set computed in the rank two subsystem. Therefore the two length
functions agree.

\subsubsection{Problem 9}

Prove that there is a unique element \(\sigma\in\mathcal W\) sending
\(\Phi^+\) to \(\Phi^-\), relative to \(\Delta\). Prove that any reduced
expression for \(\sigma\) must involve all \(\sigma_\alpha\), where
\(\alpha\in\Delta\). Discuss \(\ell(\sigma)\).

\paragraph{Hint.}
Weyl chambers are permuted simply transitively by \(\mathcal W\). Hence there
is a unique element \(w_0\) sending the fundamental chamber to the opposite
chamber, equivalently sending all positive roots to negative roots:
\[
w_0(\Phi^+)=\Phi^-.
\]
The length formula gives
\[
\ell(w_0)=|\Phi^+|.
\]
If a reduced expression for \(w_0\) omitted some simple reflection
\(\sigma_\alpha\), then \(w_0\) would lie in the parabolic subgroup generated
by \(\Delta-\{\alpha\}\). Such an element cannot send \(\alpha\) to a negative
root. Therefore every simple reflection must occur.

\subsubsection{Problem 10}

Given
\[
\Delta=\{\alpha_1,\ldots,\alpha_\ell\}
\]
in \(\Phi\), let
\[
\lambda=\sum_{i=1}^{\ell} k_i\alpha_i,
\qquad
k_i\in\mathbb Z,
\]
where either all \(k_i\geq 0\) or all \(k_i\leq 0\). Prove that either
\(\lambda\) is a multiple, possibly \(0\), of a root, or else there exists
\(\sigma\in\mathcal W\) such that
\[
\sigma\lambda=\sum_{i=1}^{\ell} k_i'\alpha_i,
\]
with some \(k_i'>0\) and some \(k_i'<0\).

\paragraph{Hint.}
Assume that \(\lambda\) is not a multiple of any root. Then the hyperplane
\[
P_\lambda=\{\mu\in E\mid (\lambda,\mu)=0\}
\]
is not contained in
\[
\bigcup_{\alpha\in\Phi}P_\alpha.
\]
Choose
\[
\mu\in P_\lambda-\bigcup_{\alpha\in\Phi}P_\alpha.
\]
Then \(\mu\) is regular, so some \(\sigma\in\mathcal W\) sends \(\mu\) into the
fundamental chamber. Thus
\[
(\alpha_i,\sigma\mu)>0
\]
for every \(i\). Since
\[
0=(\lambda,\mu)=(\sigma\lambda,\sigma\mu)
=
\sum_i k_i'(\alpha_i,\sigma\mu),
\]
and all \((\alpha_i,\sigma\mu)\) are positive, the coefficients \(k_i'\) cannot
all have the same sign. Hence some are positive and some are negative.

\subsubsection{Problem 11}

Let \(\Phi\) be irreducible. Prove that \(\Phi^\vee\) is also irreducible. If
\(\Phi\) has all roots of equal length, so does \(\Phi^\vee\), and then
\(\Phi^\vee\) is isomorphic to \(\Phi\). On the other hand, if \(\Phi\) has two
root lengths, then so does \(\Phi^\vee\); but if \(\alpha\) is long, then
\(\alpha^\vee\) is short, and vice versa. Use this fact to prove that
\(\Phi^\vee\) has a unique maximal short root, relative to the partial order
defined by \(\Delta\).

\paragraph{Hint.}
Irreducibility is equivalent to connectedness of the Dynkin diagram. Passing to
the dual root system reverses arrows but does not disconnect the diagram.
Since
\[
\alpha^\vee=\frac{2\alpha}{(\alpha,\alpha)},
\]
long roots become short coroots and short roots become long coroots. If all
roots have the same length, this operation only rescales the whole root
system. If there are two lengths, use the uniqueness of the highest long root
in \(\Phi\); its coroot gives the unique maximal short root in \(\Phi^\vee\).

\subsubsection{Problem 12}

Let \(\lambda\in C(\Delta)\). If
\[
\sigma\lambda=\lambda
\]
for some \(\sigma\in\mathcal W\), then \(\sigma=1\).

\paragraph{Hint.}
Here \(C(\Delta)\) is the open fundamental chamber, so
\[
(\lambda,\alpha)>0
\]
for every \(\alpha\in\Phi^+\). If \(\sigma\neq 1\), then its inversion set is
nonempty: there exists \(\alpha\in\Phi^+\) such that
\[
\sigma^{-1}\alpha\in\Phi^-.
\]
Then
\[
(\alpha,\sigma\lambda)
=
(\sigma^{-1}\alpha,\lambda)<0,
\]
while
\[
(\alpha,\lambda)>0.
\]
This contradicts \(\sigma\lambda=\lambda\). Hence \(\sigma=1\).

\subsubsection{Problem 13}

The only reflections in \(\mathcal W\) are those of the form
\[
\sigma_\alpha,\qquad \alpha\in\Phi.
\]
A vector in the reflecting hyperplane would, if orthogonal to no root, be fixed
only by the identity in \(\mathcal W\).

\paragraph{Hint.}
Let \(s\in\mathcal W\) be a reflection, with reflecting hyperplane \(H_s\). If
\(H_s\) were not equal to \(P_\alpha\) for any root \(\alpha\), then one could
choose a regular vector
\[
v\in H_s-\bigcup_{\alpha\in\Phi}P_\alpha.
\]
Since \(v\in H_s\), we have \(s(v)=v\). But a regular vector has trivial
stabilizer in \(\mathcal W\), by Problem 12 after moving \(v\) to the
fundamental chamber. Hence \(s=1\), contradiction. Therefore
\[
H_s=P_\alpha
\]
for some \(\alpha\in\Phi\), and the reflection is exactly \(\sigma_\alpha\).

\subsubsection{Problem 14}

Prove that each point of \(E\) is \(\mathcal W\)-conjugate to a point in the
closure of the fundamental Weyl chamber relative to a base \(\Delta\). Enlarge
the partial order on \(E\) by defining
\[
\mu<\lambda
\]
if and only if \(\lambda-\mu\) is a nonnegative \(\mathbb R\)-linear
combination of simple roots. If \(\mu\in E\), choose \(\sigma\in\mathcal W\)
for which
\[
\lambda=\sigma\mu
\]
is maximal in this partial order.

\paragraph{Hint.}
Since \(\mathcal W\mu\) is finite, choose a maximal element
\[
\lambda=\sigma\mu
\]
in its orbit with respect to this enlarged partial order. If
\[
(\lambda,\alpha)<0
\]
for some simple root \(\alpha\), then
\[
\sigma_\alpha\lambda
=
\lambda-\frac{2(\lambda,\alpha)}{(\alpha,\alpha)}\alpha
\]
is strictly larger than \(\lambda\), since the coefficient of \(\alpha\) added
is positive. This contradicts maximality. Hence
\[
(\lambda,\alpha)\geq 0
\]
for all \(\alpha\in\Delta\), so \(\lambda\) lies in the closure of the
fundamental Weyl chamber.


\section{Existence Theorem}

\subsubsection{Problem 1}

Prove that if \(\dim L<\infty\), then \(\mathfrak U(L)\) has no zero divisors.
Hint: Use the fact that the associated graded algebra
\(\operatorname{gr}\mathfrak U(L)\) is isomorphic to a polynomial algebra.

\paragraph{Hint.}
Use the standard filtration
\[
\mathfrak U(L)_0\subset \mathfrak U(L)_1\subset \mathfrak U(L)_2\subset\cdots
\]
by degree. PBW gives
\[
\operatorname{gr}\mathfrak U(L)\cong S(L),
\]
and since \(\dim L<\infty\), \(S(L)\) is a polynomial algebra, hence an
integral domain. If \(a,b\in \mathfrak U(L)\) are nonzero, take their leading
terms \(\overline a,\overline b\) in \(\operatorname{gr}\mathfrak U(L)\). Since
\(S(L)\) has no zero divisors,
\[
\overline a\,\overline b\neq 0,
\]
so the leading term of \(ab\) is nonzero. Hence \(ab\neq 0\).

\subsubsection{Problem 2}

Let \(L\) be the two dimensional nonabelian Lie algebra (1.4), with
\[
[x,y]=x.
\]
Prove directly that
\[
\iota:L\to \mathfrak U(L)
\]
is injective, i.e.
\[
\ker\iota\cap L=0.
\]

\paragraph{Hint.}
Use the defining relation in \(\mathfrak U(L)\):
\[
xy-yx=x.
\]
One direct way is to construct a representation of \(L\). Take
\[
X=E_{12},\qquad
Y=
\begin{pmatrix}
-1&0\\
0&0
\end{pmatrix}.
\]
Then
\[
[X,Y]=X.
\]
By the universal property of \(\mathfrak U(L)\), this gives an algebra
homomorphism
\[
\mathfrak U(L)\to \operatorname{End}(F^2)
\]
sending \(x\mapsto X\), \(y\mapsto Y\). If \(ax+by\) maps to zero in
\(\mathfrak U(L)\), then applying this representation gives
\[
aX+bY=0.
\]
Since \(X,Y\) are linearly independent matrices, \(a=b=0\). Hence
\(\iota\) is injective on \(L\).

\subsubsection{Problem 3}

If \(x\in L\), extend \(\operatorname{ad}x\) to an endomorphism of
\(\mathfrak U(L)\) by defining
\[
\operatorname{ad}x(y)=xy-yx,
\qquad y\in\mathfrak U(L).
\]
If \(\dim L<\infty\), prove that each element of \(\mathfrak U(L)\) lies in a
finite dimensional \(L\)-submodule. If \(x,x_1,\ldots,x_m\in L\), verify that
\[
\operatorname{ad}x(x_1\cdots x_m)
=
\sum_{i=1}^m x_1\cdots x_{i-1}(\operatorname{ad}x)(x_i)x_{i+1}\cdots x_m .
\]

\paragraph{Hint.}
First show that \(\operatorname{ad}x\) is a derivation of the associative
algebra \(\mathfrak U(L)\):
\[
\operatorname{ad}x(uv)=xu v-uvx=(xu-ux)v+u(xv-vx).
\]
This gives the displayed formula by induction on \(m\). Now use the PBW
filtration. If \(u\in \mathfrak U(L)\) has degree at most \(m\), then
\(\operatorname{ad}x(u)\) also has degree at most \(m\). Since \(\dim L<\infty\),
the filtered piece \(\mathfrak U(L)_m\) is finite dimensional. Hence the
\(L\)-submodule generated by \(u\) is contained in \(\mathfrak U(L)_m\), and is
therefore finite dimensional.

\subsubsection{Problem 4}

If \(L\) is a free Lie algebra on a set \(X\), prove that \(\mathfrak U(L)\) is
isomorphic to the tensor algebra on a vector space having \(X\) as basis.

\paragraph{Hint.}
Let \(V\) be the vector space with basis \(X\). The tensor algebra \(T(V)\) has
the universal property that any map
\[
X\to A
\]
to an associative algebra \(A\) extends uniquely to an algebra homomorphism
\[
T(V)\to A.
\]
On the other hand, since \(L\) is the free Lie algebra on \(X\), any map
\[
X\to A_{\operatorname{Lie}}
\]
extends uniquely to a Lie homomorphism
\[
L\to A_{\operatorname{Lie}}.
\]
Then the universal property of \(\mathfrak U(L)\) extends this uniquely to an
algebra homomorphism
\[
\mathfrak U(L)\to A.
\]
Thus \(\mathfrak U(L)\) and \(T(V)\) have the same universal property, so they
are naturally isomorphic:
\[
\mathfrak U(L)\cong T(V).
\]

\subsubsection{Problem 5}

Describe the free Lie algebra on a set
\[
X=\{x\}.
\]

\paragraph{Hint.}
With only one generator, every bracket involving only \(x\) is zero because
\[
[x,x]=0.
\]
Thus no new nonzero elements are produced by taking Lie brackets. Hence the
free Lie algebra on one generator is just the one-dimensional abelian Lie
algebra
\[
Fx.
\]
Check this using the universal property: any map from \(\{x\}\) to a Lie
algebra \(M\) extends uniquely to a Lie homomorphism
\[
Fx\to M.
\]

\subsubsection{Problem 6}

How is the PBW Theorem used in the construction of free Lie algebras?

\paragraph{Hint.}
PBW is used to control the embedding of a Lie algebra into its universal
enveloping algebra. In particular, PBW implies that the canonical map
\[
L\to \mathfrak U(L)
\]
is injective. For a free Lie algebra generated by \(X\), Exercise 4 gives
\[
\mathfrak U(L)\cong T(V),
\]
where \(V\) has basis \(X\). Thus \(L\) can be realized as the Lie subalgebra
of the tensor algebra \(T(V)\) generated by \(X\), where the Lie bracket is the
commutator
\[
[a,b]=ab-ba.
\]
PBW guarantees that this realization does not collapse \(L\); equivalently, the
Lie words in \(X\) satisfy only the necessary Lie algebra identities.
\end{document}